% Standalone sections for inclusion in the handoff paper. No preamble.
% Requires only the standard mathematical commands supplied by amsmath/amssymb.

\section{Scope of the all-age and full-history argument}
\label{ha:scope}

The purpose of these sections is to separate finite algebra, positive
majorant estimates, and functional-analytic consequences from the analytic
source attachment on which their application depends.  The resulting
statement is conditional.  In particular, the estimates below do not
establish a nondegenerate twenty-one-component field, its continuous tensor
transformation law, or a full Yang--Mills theory.

Fix a finite regulator, one source-free trajectory, and a terminal weak
coupling band.  Primitive dates are denoted by $j\leq N$, physical cell
lengths by $r_j=\ell^{j-N}$, and inverse couplings by
$\gamma_j=g_j^{-2}$.  We distinguish the physical analytic radius
$b_j=cg_j\log\gamma_j$ from the original source radius
$a_j=Cg_j\log^{q_{\rm rad}}\gamma_j$.  The retained source fluctuation
gain is
\begin{equation}
 \delta_j=\frac{g_jT_j}{a_j},\qquad
 \delta_j\leq C(\log\gamma_j)^{-\nu},\quad \nu>0.
 \label{ha:delta}
\end{equation}
It must not be identified with the physical regular-graph parameter
$\eta_j=(g_j/b_j)^2$.  Aligned endpoints have positive drift; fixed
primitive oscillations are included in constants.  The estimates used
below include
\begin{equation}
 \gamma_i\leq\gamma_j+D(1+j-i),\qquad
 j-i\leq C\gamma_i,
 \label{ha:drift}
\end{equation}
on the weak part of the trajectory.  The terminal band is fixed before
taking ultraviolet limits.

An original activity is retained with its birth date, pointed root,
carrier, original independent-pair extension, and scalar coordinate.
An original chronological operator is estimated on every admitted analytic
test function.  A bound on its value at $1$, or a normalized expression
$T_H/T_H1$, is never substituted for that operator norm.  A regular
physical field norm is not an original independent-$(U,J)$ norm.

Here are the principal application hypotheses used in this part of the
paper.  Their distinction is important.
\begin{enumerate}
\item The actual regular source difference maps have the stated
same-contour analytic margins, complete supported norms, fixed-degree
marked cluster bounds, and positive dated creation majorants.  Their
physical marginal graph and its future-corrected projection are available
on the physical $b_j$ balls.
\item The original chronological expansion has a convergent positive
rooted majorant, its all-spectator composition norm, and fixed spare
occupation, carrier, and component-interior reserves.  All reference
inverses are the original source-free ones.
\item At a first mixed-history read, the actual nested minimum-plus-response
maps have the required local regularity on the complete averaging and
current footprints, on the same localization contours as their neutral
bound.  This is the unresolved source-attachment hypothesis discussed
below; final-field admission alone does not imply it.
\item When two global source colours are used, carried mixed coefficients
have the retained two-slot profile representation and original spectator
majorant stated in Section~\ref{ha:two-colour}.  Prefusion live seeds
retain their unused early reserve.  These are stronger assertions than
two separate first-mark bounds.
\end{enumerate}
Statements explicitly described as algebraic do not require the third
hypothesis.  The copied estimate and its downstream field construction do.

\paragraph{Provenance and status precedence.}
The status controlling these sections is
\path{work/ym4-resumed-research-status-2026-10-18.txt}.
Earlier statements that a repeated integration audit had closed every
analytic input are superseded by the admission gap in
\path{work/static_original_tube_first_contact_audit_2026-10-18.txt},
Sections 11--15.  Conversely, the later fixed-coupling same-state
reflection argument discussed next supersedes the earlier torus-topology
objection; these two updates concern different issues.

\section{Scalar prerequisites and their precise logical role}
\label{ha:scalar}

The later operator construction takes place on one scalar Hilbert space.
It therefore requires more than the formal existence of scalar
off-diagonal correlations.  We record the scalar inputs and their
conditional consequences rather than assuming a scalar Poincare theory
without explanation.

\subsection{Flat-diagonal growth and the reflection space}

Let $\mathcal D_n=\bigcup_{i<j}\{x_i=x_j\}$ and let
$\mathcal S^0_n$ be the Schwartz functions flat, to every order, on
$\mathcal D_n$.  Suppose compatible off-diagonal correlations obey, on
separated tensor packets,
\begin{equation}
 |S_n(f_1\otimes\cdots\otimes f_n)|
 \leq A^n(n!)^B(1+d^{-1})^{cn}
               \prod_{i=1}^n\|f_i\|_{C^1},
 \label{ha:packet-growth}
\end{equation}
where packet radii are at most $\rho/n$, their mutual separation is $d$,
and the constants are independent of placement.  There are fixed
$A_1,B_1,C_1$ such that the canonical extension satisfies
\begin{equation}
 |S_n(F)|\leq A_1^n(n!)^{B_1}p_{\lceil C_1n\rceil}(F),
 \qquad F\in\mathcal S^0_n.
 \label{ha:flat-growth}
\end{equation}
Here $p_m(F)=\max_{|\alpha|\leq m}\sup_x
(1+|x|)^m|\partial^\alpha F(x)|$; changing equivalent norms changes
the constants only.

\emph{Proof of the extension mechanism.}
Use a capped dyadic configuration partition of mesh
$\varepsilon_s=2^{-s}c/n$, supported where every pair is separated
by a fixed multiple of $n\varepsilon_s$.  It can be formed from
$O(n^2)$ fixed Gevrey transition factors and a product spatial partition.
At total derivative order $O(n)$, the product rule and Gevrey bounds cost
$n^{O(n)}$, rather than an uncontrolled constant exponential in $n^2$.
Rescale each configuration box and expand in a tensor Fourier series.
Six derivatives in each four-dimensional variable make the weighted
Fourier series summable after the $C^1$ cost in
\eqref{ha:packet-growth}.  A box then costs
$\varepsilon_s^{-(c+1)n}$ times derivatives of $F$ of order at most
$6n$.  Flatness supplies $\varepsilon_s^{kn}$ by Taylor expansion toward
a nearest full collision, with $k>c+5$.  Summing configuration boxes
costs $\varepsilon_s^{-4n}$, leaving the summable power
$\varepsilon_s^{(k-c-5)n}$.  Spatial Schwartz weights sum with order
linear in $n$.  Every remaining $n^{Dn}$ is bounded by
$e^{Dn}(n!)^D$.  This proves \eqref{ha:flat-growth}.  Compact tests
away from $\mathcal D_n$ are dense in $\mathcal S^0_n$: first cut off
at infinity, then remove a collision neighborhood, using arbitrarily
high flatness to absorb each fixed derivative of the cutoff.  This
also proves uniqueness.  If the regulator estimates require the mesh
to be smaller than the packet, first take limits on fixed separated
packets and only then perform this extension.  No uniform estimate
below the regulator mesh has been inferred. \hfill$\square$

For a mirror with normal $e$, let $\mathcal D_e$ consist of finite test
sequences whose nonconstant components are smooth, compact, internally
off full collisions, and supported where all $e\cdot x_i>0$.  Define
\begin{equation}
 (F,G)_e=\sum_{n,m}S_{n+m}
       ((\theta_e\overline{F_n})^{\rm rev}\otimes G_m).
 \label{ha:rp-form}
\end{equation}
Finite-regulator reflection positivity passes to these tests whenever
the same-state limiting construction supplies that mirror.  Complete
the quotient by its null space.  This is a Hilbert-space construction
from distributions, not an ordinary Euclidean probability $L^2$ space.
The test space is a vector space and a partial separated-support algebra;
arbitrary products can create collisions and are not automatically in it.

Positive normal translations $T(t)$ satisfy
\begin{equation}
 \|T(t)F\|\leq \|F\|^{1-2^{-m}}
                         \|T(2^mt)F\|^{2^{-m}}.
 \label{ha:semigroup-doubling}
\end{equation}
The rightmost norm grows polynomially in its translation argument by
Schwartz continuity.  Sending $m\to\infty$ proves contractivity, and
test continuity gives a strongly continuous semigroup $e^{-tH}$ with
$H\geq0$.  Tangential translations give strongly continuous unitaries.
No cutoff of equal-time hypersurfaces is used: such hypersurfaces contain
points off the full collision diagonal and cannot generally be removed
in the original Schwartz topology.

\subsection{Same-state reflections, slabs, and reconstruction}

For the literal Wilson $SU(2)$ measure at each fixed finite bare coupling,
the audited two-frame argument yields diagonal reflection positivity of
each coordinate-periodic thermodynamic subsequential state.  Its content
is stronger than equality of pressures.  Adapted diagonal tori first
give raw-link, closed-halfspace positivity.  Twisted transfer traces are
bounded by Schatten Holder inequalities; their straightened factors
are ordinary Wilson integrals.  Sparse tilts and operator-norm
amplification produce a comparison law.  An exact paired local reduction
then yields coordinate reflection invariance and even translations of
that law.  Convex separation, with gauge averaging of a full-space
separator, puts the original coordinate state in the full diagonal-RP
class.  The amplification is performed at fixed coupling, volume tends
to infinity before the tilt tends to zero, and finite-rank approximation
is taken after those bounds.  This argument does not assert RP on every
straight finite torus, and it is not itself a continuum rotation theorem.

This volume-first statement is compatible with the stopping regulators.
For fixed bare $A$, bounded positive block drift bounds the stopping
date by a finite $n_{\max}(A)$.  Take arbitrarily large compatible tori
supporting every step up to that date, and pass to a volume subsequence
on which the stopping integer $N_A$ is constant.  The resulting law is
a fixed-$A$ coordinate-periodic thermodynamic limit.  Sending
$A\to\infty$ afterward gives $N_A\to\infty$.  Additional countable
subsequences can fix the finite witness offsets.  This does not require
finite-volume running coefficients to be identical at different volumes.

The scalar slab/rotation route still requires its source pressure and
identification inputs.  In a coordinate plane, a normalized packet
insertion with two positive paddings has the bound
$C\delta^{-7/2}$.  Reflection doubling, with endpoint factors assigned
outside the array cell, turns the source-array estimate into this
operator bound.  A planar cone comparison using the coordinate and
diagonal RP spaces gives $H\geq |P_i|$ in each selected plane.
For example its two spectral representations imply
\begin{equation}
 \int e^{-2L\omega}|F(\omega-ip)|^2\,d\rho
 \leq C_A\sup_{r\geq0}|F(r)|\sup_{s\in\mathbb R}|F(is)|.
 \label{ha:cone-probe}
\end{equation}
The Laplace probes
$F_j(z)=e^{-j}z^{1-4i\epsilon j/\pi}(1+z)^{-3}$
have bounded product of the two suprema and grow exponentially on
$\epsilon\arg z>\pi/4$.  Positivity of $\rho$ excludes $|p|>\omega$.
No isometric identification of the two RP Hilbert spaces is required.

The bounded insertion construction supplies a \emph{joint} tube and
mixed Gram identities.  Separate one-gap analyticity would not suffice.
With the same-state two-frame input, the angular argument and the
$7/2<4$ bound give the scalar planar rotation continuation.  A full-point
version has heat bound $C\tau^{-5}$ and off-collision kernel estimate
$C^n n^{10n}(1+d^{-5n})$.  Its covariant extension of the ordered sector,
if these identifications are retained, obeys
\eqref{ha:flat-growth}, symmetry, and RP.  It need not agree with an
arbitrary earlier extension on equal-time contact configurations.

The corrected Osterwalder--Schrader reconstruction theorem on
$\mathcal S^0$ consumes linear seminorm order and a fixed factorial
power, as in \eqref{ha:flat-growth}; it does not require that power to
be small.  With Euclidean covariance, RP, and symmetry it gives the
scalar Wightman theory, leaving clustering as a separate input.
Thus the scalar theory assumed later is conditional on the actual
source estimates and the joint tube/ordered-sector identification,
not on a choice of arbitrary full-Schwartz contact extensions.

The full-point insertion bound also identifies
$e^{-\tau H}\phi(f)e^{-\tau H}$ with a bounded Minkowski-smeared
operator of norm $C\tau^{-5}\|f\|_1$.  For compact energy-transfer
support, unit energy bands have bounded width of interaction.  Taking
$\tau=(1+n+m)^{-1}$ gives block norm
$C\|f\|_1(1+n+m)^5$; block Schur summation proves
$\|\overline{\phi(f)}(1+H)^{-5}\|<\infty$.  The identification uses
the joint tube and common test domain before removing external imaginary
gaps.  A bounded Euclidean insertion alone is not this identification.
The generalized Hamiltonian bound permits the Driessler--Summers
decomposition into clustering scalar Wightman phases with unique vacua,
under that theorem's hypotheses.  It is not generic ergodic decomposition
of an assumed Euclidean probability law.

A nontrivial phase can be retained with a derivative witness
$f=\partial_t\partial_{x_1}\psi$, supported at positive Euclidean time.
The exact strong test-translation derivative gives
$v_f=-iHP_1v_\psi$, with the appropriate dilation factor for dilated
packets.  A strictly positive reflected norm therefore has a nonzero
projection into some compact positive-energy interval, outside
$\ker P_1$.  This projection disintegrates into nonnegative phase norms,
so a positive-measure collection of the decomposed phases is nontrivial.
Neither a spectral mass gap nor phasewise ultraviolet normalization
follows.  An averaged asymptotic coefficient can vary between phases;
phasewise matching requires a decomposable operator estimate or estimates
with all bounded central insertions.

\paragraph{Provenance.}
The detailed scalar inputs are recorded in
\path{work/os_s0_whitney_extension_2026-10-18.txt},
\path{work/s0_reflection_common_core_audit_2026-10-18.txt},
\path{work/same_state_diagonal_rp_reaudit_2026-10-18.txt},
\path{work/rotation_slab_operator_consumer_2026-10-18.txt},
and \path{work/scalar_reconstruction_and_phase_selection_2026-10-18.txt}.
The opening diagonal-RP warning in the older common-core note is superseded
by the later fixed-coupling re-audit, not by an automatic continuum claim.

\section{Permanently dated periodic marginals}
\label{ha:periodic}

Maturation of a complete cohort does not erase the individual old
phase-root marginals.  The following repair is an identity on the birth
root lattice and on the same original fine field.

\paragraph{Periodic primitive.}
Let $a:\mathbb Z^4\to V$ be $p$-periodic with zero mean, where $V$ is
a normed tensor space or finite coefficient algebra.  Write $A_\mu$ for
averaging in coordinate $\mu$ and set
$f_\mu=A_1\cdots A_{\mu-1}a-A_1\cdots A_\mu a$.
Each $f_\mu$ has zero sum on its $\mu$-line.  Its cumulative sum, extended
periodically, gives
\begin{equation}
 a(z)=\sum_{\mu=1}^4D^-_\mu b_\mu(z),\qquad
 \|b_\mu\|_\infty\leq2p\|a\|_\infty,
 \quad
 \sum_z a(z)F_z=-\sum_{z,\mu}b_\mu(z)D^+_\mu F_z.
 \label{ha:coboundary}
\end{equation}
The wrap-around equality follows from the zero line sum.  Rectangular
periods use their maximum.  The construction is linear, so label sums
may remain inside the spatial norm throughout.

Let $q_{AB}(U,x)$ be the same fine clover tensor, with its fine volume
factor, at every date.  Set
\begin{equation}
 \widehat Q_{i,z}(U,J)=\sum_xh^i_z(x)q_{AB}(U,x),\qquad
 h^i_z(x)=\prod_\mu(1-|x_\mu/r_i-z_\mu|)_+.
 \label{ha:qhat}
\end{equation}
This reference is independent of the formal current $J$.  A change of
fine summation index gives the exact identity
\begin{equation}
 \widehat Q_{i,z+e_\mu}-\widehat Q_{i,z}
 =\sum_xh^i_z(x)\{q_{AB}(U,x+r_ie_\mu)-q_{AB}(U,x)\}.
 \label{ha:qhat-difference}
\end{equation}
For an actual eligibility mask $\chi$ and a profile $f$, summation by
parts in \eqref{ha:coboundary} differentiates $\chi(z)f(r_iz)\widehat Q_{i,z}$.
Its exact product rule has three terms: the fine-clover difference,
the profile difference, and the mask difference.  The latter two keep
their original endpoints and birth widths.  They are not discarded
as changes of representation.

On a common admitted smooth read of physical scale $r_n$, put $h=r_i/r_n$.
The same-contour curvature and Holder estimates imply
\begin{equation}
 |\widehat Q_{i,z}|\leq Ca_n^2h^4,
 \qquad |D^+\widehat Q_{i,z}|\leq Ca_n^2h^{4+\beta}.
 \label{ha:phase-gain}
\end{equation}
There are $O(h^{-4})$ old roots per $n$-cell.  The phase sum is therefore
bounded by $Cpa_n^2\|a\|_\infty h^\beta$, whereas its direct neutral
bound is $Ca_n^2\|a\|_\infty$ with \emph{no period factor}.  Taking
the minimum before the age sum is decisive.  A profile difference costs
$pr_i^\theta[f]_\theta$; an actual boundary commutator is supported in
an old-width strip, with any printed $MR_i$ collar retained.

At an artificial boundary only a patch of scale $d r_n$ may be smooth.
The safe functional comparison is the minimum of a neutral and a
prepared Ward bound, of the form
\begin{equation}
 Ca_n^2h^4\min\{1,Kp h^\beta d^{-(4+\beta)}\}.
 \label{ha:graded-phase}
\end{equation}
The same function and contours are required on both sides.  Summing
dyadic collar shells, whose relative volume is $O(d)$, with
$t=(Kph^\beta)^{1/(4+\beta)}$, costs $O(t)$: below $t$ use the neutral
bound; above $t$ sum $Kph^\beta d^{-(3+\beta)}$.  This proves a positive
gain of exponent $\beta/(4+\beta)$, capped at one, without asserting
uniform smoothness through the seam.  The use of this estimate for the
actual mixed maps is precisely part of the unresolved attachment.

\paragraph{Exact matured mean.}
If $R=r_j/r_i$ is divisible by $p$, then for every residue vector $\alpha$,
\begin{equation}
 R^{-4}\sum_{z\equiv\alpha\ ({\rm mod}\ p)}
       \prod_\mu(1-|z_\mu/R-Z_\mu|)_+=p^{-4}.
 \label{ha:tent-balance}
\end{equation}
In one dimension write $R=pq$ and partition the translates of the tent
into $q$ unit partition-of-unity families; the sum is $q$.  Products
prove the formula.  Thus the grouped affine curvature coefficient of
a periodic old reference is exactly its phase mean.  Its oscillatory
full function still exists and is bounded by
\eqref{ha:coboundary}--\eqref{ha:graded-phase}.

Split each complete pointed birth at creation as
$v_i+\widehat Q_i d_i+Z_i$, with its full affine curvature projection
$\Pi Z_i=0$.  Define $M_i$ to bound the \emph{whole} nonconstant newborn,
including $d_i$, not just $Z_i$.  Exact affine reproduction gives
\begin{equation}
 \Pi_jW_j(Z)=\sum_{i\leq j}R_{ij}^{-4}
                 \sum_zh_Z(z/R_{ij})d_i(z).
 \label{ha:physical-marginal}
\end{equation}
Old zero-projection remainders contribute zero because their first jet
vanishes and the linear prefix reproduces affine curvature.  Their
quantum changes are already later births.  Mature terms in
\eqref{ha:physical-marginal} are phase means; the remaining positive hat
averages are bounded by $C\sum_{i\,{\rm young}}M_i$.  Youth lasts only
$O(\log p_i)$ dates.  No old function is redated, and no scalar or
historical boundary term is removed by this identity.

\paragraph{Provenance.}
\path{work/periodic_marginal_coboundary_2026-10-18.txt};
the hats and pointed decomposition are those associated with CMP119
(3.40)--(3.43), (3.65)--(3.67), with tensor projection before its scalar
Euclidean reduction.

\section{Positive Volterra estimates and a separate native-error lane}
\label{ha:volterra}

Suppose the literal dated preparation kernel obeys, for $t=j-i$,
\begin{equation}
 k_{ji}\leq C(1+t)^q
 \min\{1,C(1+\log\gamma_i)^r e^{-\beta t}\}
             +C(1+t)^q e^{-\beta t}.
 \label{ha:phase-kernel}
\end{equation}
Put $L_j=1+\log(1+\log\gamma_j)$.  Then
\begin{equation}
 \sum_{i\leq j}k_{ji}\leq C L_j^{q+1}.
 \label{ha:row-sum}
\end{equation}
Indeed, for $t\leq\gamma_j$, \eqref{ha:drift} makes $\gamma_i$ comparable
to $\gamma_j$.  Split at $C_0+C_1\log\log\gamma_j$: use the neutral
term below and the exponential above.  For $t>\gamma_j$, use
$\gamma_i\leq C(1+t)$ and an absolutely summable exponential tail.
Therefore $\delta_jL_j^{q+1}\to0$ at a sufficiently weak floor.

If a matured physical cohort has
$|c_j|\leq A(\gamma_i/\gamma_j)^p$ and its actual newborns satisfy
\begin{equation}
 M_{j+1}\leq C\delta_j\left(|c_j|+
                  \sum_{u=i}^jk_{ju}M_u\right)+f_{j+1},
 \label{ha:birth-rec}
\end{equation}
with $M_i\leq A$ and $f_{j+1}\leq A(\gamma_i/\gamma_{j+1})^p$, then
\begin{equation}
 M_j\leq CA(\gamma_i/\gamma_j)^p,\qquad
 \sum_{u=i}^jM_u\leq CA\gamma_i(\gamma_i/\gamma_j)^p.
 \label{ha:birth-poly}
\end{equation}
Divide by the monotone weight $(\gamma_i/\gamma_j)^p$ and use
\eqref{ha:row-sum}; choosing the feedback coefficient below $1/2$
proves the first bound by induction.  The number of dates in
\eqref{ha:drift} proves the second.  This is a positive bound for a
list of original births, not a static semigroup on arbitrary pair tubes.

On a fusion-to-midpoint interval where $\gamma_j\asymp G$, the same
proof with weight $e^{\alpha(j-i)}$, $\alpha=c/(1+\log\log G)$,
gives
\begin{equation}
 M_j\leq CA\left(e^{-c(j-i)/(1+\log\log G)}+\zeta_G\right)
 \label{ha:slow-decay}
\end{equation}
when the physical forcing has a stable exponential plus the floor
$\zeta_G$, and the added forcing is $O(A\delta_j\zeta_G)$.
The weighted row sum is unchanged up to a constant: the neutral
memory interval has length $O(\log\log G)$ and the tail still decays.
Across $cG$ dates the first term beats every fixed
$\exp[C(\log G)^d]$ preparation loss.  A power floor $G^{-K}$ does not.

\subsection{The actual old/young split}

Let $\rho_j$ be the actual regular density, $Y_j$ its source-free young
$R$ sum, and $\sigma_j=\rho_je^{-Y_j}$.  The regular comparison is
\begin{equation}
 C_j^0W=\frac{K_j^{\rm in}(\sigma_jW)}{K_j^{\rm in}\sigma_j},
 \qquad C_jW=\frac{K_j^{\rm in}(\rho_jW)}{K_j^{\rm in}\rho_j}.
 \label{ha:actual-core}
\end{equation}
The first kernel contains the old $R$ whose partitions fit the next
lattice.  It is neither an all-$R$ kernel nor an $E$-only kernel.
Its successive densities satisfy the exact identity
\begin{equation}
 \sigma_{j+1}=e^{A_j}K_j^{\rm in}\sigma_j,\quad
 A_j=F_j+h_j+log\frac{K_j^{\rm in}\rho_j}
                              {K_j^{\rm in}\sigma_j}-Y_{j+1}.
 \label{ha:intermediate-normalizer}
\end{equation}
The intermediate $A_j$ remain in a joint block integral.  Dropping them
would change the product of normalized kernels.

Interpolating only the source-free $Y_j$ gives, with $P_j$ the unchanged
static prefix,
\begin{equation}
 (C_j-C_j^0)W=\int_0^1
       \operatorname{Cov}_{\sigma_je^{sY_j}}
       (W-W(P_j),Y_j-Y_j(P_j))\,ds.
 \label{ha:native-identity}
\end{equation}
This is exact, annihilates constants, and acts on the same complete mark.
It does not bound derivatives of a sourced young-$R$ family by a native
source-free power.  The original young renewal estimates yield a literal
newborn bound, including its scalar value,
\begin{equation}
 a_{j+1}^{-2}\|(C_j-C_j^0)W\|_{\rm new,orig}
 \leq C\gamma_j^{d}g_j^{K_{\rm nat}}
                     a_j^{-2}\|\Delta W\|_{\rm prep}.
 \label{ha:native-newborn}
\end{equation}
The powers in this estimate are fixed before the native target exponent.
Young parent ages are $O(\log\log\gamma_j)$, so
$(g_i/g_j)^{K_{\rm nat}}\leq
\exp[CK_{\rm nat}\log\log\gamma_j/\gamma_j]$.
Thus stronger native output powers can be obtained by increasing the
weak floor, while the admission exponent $\kappa_R$ remains fixed.
Fresh healed $R'$ uses its genuine bad-field reserve; regular $R''$
renewal uses \eqref{ha:native-identity}, not that reserve.

Unroll a difference by its first use of $C-C^0$.  Early first-native
outputs are already small in their original creation norm and continue
through the polynomial cohort bound \eqref{ha:birth-poly}.  Late
first-native outputs act on the already decayed prepared difference.
Consequently a native power-small error is never multiplied by the
large graded-read constant.  Choosing $K_{\rm nat}$ after the fixed
admission, creation, transport, and date-count powers leaves any desired
fixed power $G^{-K}$.  The choice is acyclic only if strengthening the
native estimate changes the weak floor and constants, not those powers.

\paragraph{Provenance.}
\path{work/phase_kernel_positive_volterra_2026-10-18.txt},
\path{work/singleton_dated_birth_polynomial_operator_2026-10-18.txt},
\path{work/native_R_power_vs_source_admission_2026-10-18.txt},
and \path{work/native_young_R_marked_error_lane_2026-10-18.txt}.

\section{Causal scalar extraction and complete-copy coefficients}
\label{ha:scalars}

Introduce a nilpotent variable $\xi_a$, $\xi_a^2=0$, only after copy $a$
has been assembled as a complete scalar density coefficient, including
all its lower subsets, deterministic matching coefficients, and identity
term.  A copy is not one of the partial blocks of its logarithm.
In a local norm the sum over label sets is inside the spatial sum:
\begin{equation}
 \|b\|_{\mathbf t,\kappa}=
 \sup_x\sum_{H\ni x}e^{\kappa d(H)}
                \sum_{S\ne\varnothing}t_S\|b_{H,S}\|,
 \qquad t_S=\prod_{a\in S}t_a.
 \label{ha:local-algebra}
\end{equation}
A sum of global norms of separate copies would grow with their number
even when their supports are disjoint.

At a universal birth $l$, write the actual coefficient
$A_{l,\alpha}=v_{l,\alpha}+B_{l,\alpha}$, where
$B_{l,\alpha}(I,0)=0$.  Let $I_l(H,\alpha)$ be its actual birth eligibility.
Then, identically in the finite algebra,
\begin{equation}
 \exp\!\left(\sum_{I_l=1}A_{l,\alpha}\right)
 =e^{V_l}e^{-D_l(H)}
       \exp\!\left(\sum_{I_l=1}B_{l,\alpha}\right),
 \quad
 V_l=\sum_\alpha v_{l,\alpha},\quad
 D_l=\sum_{I_l=0}v_{l,\alpha}.
 \label{ha:birth-scalar}
\end{equation}
The central $V_l$ commutes with every later frozen linear operator.
$D_l$ retains its birth history and witness.  For any later mask $J$,
the replacement of $I_l$ by $J$ must retain the opposite scalar mask
correction before taking absolute values.  A later hole cannot create
a new microscopic deficit from an already exported scalar.  The
identity applies to genuine universal healed births as well as regular
births; a history-dependent boundary scalar was never universal and
remains in its chronological operator.

This order has a quantitative consequence.  At birth,
$\Omega_l\supset\Lambda_l\supset\Omega_{l+1}$ implies
$\Lambda_l^c\subset\bigcup_{n\leq l}\Gamma_n$.  Failure of the literal
pointed restrictions $z\in\Lambda_l^{(0)}$, $X\subset\Lambda_l$ gives
either the prescribed current collar or a retained carrier path to that
complement.  Per actual own-$n$ cell, $n\leq l$, the number of current-$l$
roots at carrier distance $q$ is at most
$C(1+q+MR_l)^4$, since $r_n/r_l\leq1$.  A tree reserve pays the $q$
moment.  There is no future scale $p$ and no factor $(r_p/r_l)^4$.
This census concerns actual new scalar coefficients.  An old fusion-root
coefficient cannot be relabelled as a new birth to use it.

For an opened component, suppose the assigned singleton deficits satisfy
$|d_a(H)|\leq P(G)n_a(H)$ and $\sum_an_a(H)\leq CD_H$ in its actual
paid cell/path census.  Nilpotence gives
\begin{equation}
 \left\|e^{\sum_a\xi_ad_a(H)}\right\|_{\mathbf t}
 =\prod_a(1+t_a|d_a|)
 \leq \exp\!\left(\sum_an_a\log(1+t_aP(G))\right).
 \label{ha:deficit-product}
\end{equation}
Choosing the logarithm below a reserved fraction of the actual occupation
debit pays this factor on every analytic spectator.  The inequality is
applied within the labelled branch, before the history sum.  It does not
tilt the unmarked measure by occupation of every historical bad cell.
Nor does exact full-state centering prove this estimate.  Universal
singleton scalar sums may be extensive and stay outside it.

\paragraph{Provenance.}
\path{work/stopped_scalar_chronological_normal_form_2026-10-18.txt},
\path{work/initial_clone_census_quotient_audit_2026-10-18.txt},
and \path{work/copy_pressure_normalization_integration_audit_2026-10-18.txt}.

\section{The first-contact classification and the remaining source gap}
\label{ha:first-contact}

Consider one complete adjacent-frame difference whose future-corrected
regular coordinate vanishes.  Choose a midpoint $m$ with
$\gamma_m\asymp G$ and $m-f\geq cG$; lower it by the finite preparation
window where necessary.  At each old root keep its full original
carrier, clover, current stencil, and attachments.  The exhaustive
estimate distinguishes the following cases.

An early genuine bad-field or annular event retains its original birth
and attachment reserve. If an old short stencil first meets an ancient
raw cell at a later read, and the cell tier $n$ is no finer than the
birth tier $i$, the original neutral fourth power cancels the old-root
census at that own tier. For $n<i$, retain the literal original-domain
operator admission and the actual cell/path occupation; one must not
rescale the birth into the finer tier and claim that cancellation.
A root far from its copy instead keeps
its original exit path.  Remote dependence on a bad cell does not itself
count as a nearby contact.  If the short stencil meets no old actual cell,
its existing determining tier is late, apart from artificial graded seams
and the preparation cap; the improved read is applied there.  Newly
introduced late roots have their own small creation norm and do not need
a fabricated early path.  Periodic, mask, and profile terms retain the
identities in Sections~\ref{ha:periodic}--\ref{ha:scalars}.

For the smooth case the target estimate is a same-contour minimum of
the original neutral bound and a proper-layer Ward bound.  A graded
collar sum and the phase recurrence then give, schematically,
\begin{equation}
 P(G)L_k\exp\!\left[-\frac{cG}{1+\log\log G}\right],
 \qquad
 L_k\leq\exp[C(1+R_k)(1+\log\gamma_k)^C].
 \label{ha:proper-layer-gain}
\end{equation}
This beats the fixed preparation loss.  The native lane of
\eqref{ha:native-newborn} is kept separate.  After all fixed polynomial
losses the intended complete-copy small factor is
\begin{equation}
 \varepsilon(G)=e^{-c(\log G)^s}+G^{-K},\qquad s>1.
 \label{ha:epsilon}
\end{equation}

\paragraph{Unresolved application hypothesis.}
The estimate \eqref{ha:proper-layer-gain} is not supplied by membership
of the final field in the original independent-pair domain.  One needs
the actual nested bases
$U(\mathcal B(Y_r);\hbox{modified coarse data})$, together with their
independent-current companions and localized responses, to be admitted
on the complete local averaging/operator footprint in one compatible
gauge.  Both neutral and improved bounds must hold on the same contour.
The local current inverse now gives a physical current estimate on an
\emph{already admitted} finite same-family collar.  It does not create
that collar across a nested minimum.

The latest slotwise reduction references the preceding complete
reconstructed field.  New boundary slots then agree exactly, and old
interior slots retained by its baseline differ by controlled response
and the retained endpoint gauge.  The unresolved slots are the
preparation-changed ones, including their transition back to retained
blocks, on their full footprints.  Independently good gauges on two
pieces do not give one admissible mixed coarse datum.  Peeling an outer
response does not invert a nested minimum.  Deleting an entire coarse
footprint as a neutral strip leaves a cost of order $r_p/r_n$, which
need not tend to zero with $G$.  A graded bypass would require its
terminal collars to remain in one admitted branch or a common identity
annulus; the required inclusion is not verified.

A dimension-independent real-slice continuation can weaken the complex
target: a bound on a full real parameter diameter and a fixed complex
margin yields a positive reduced exponent by the two-constants theorem.
A physical graph or one real specialization is not that full diameter.
Thus this alternative does not remove the present admission hypothesis.
The subsequent assembly is a theorem from that hypothesis and the
specified original majorants, not an independent proof of it.

\section{Merged groups and a single prefix-tree reserve}
\label{ha:forest}

Copy supports, including their read collars, are separated by $\tau$.
For one term of the original connected expansion retain its paid
component interiors and attachments, with their dates.  Let $T_i$ be
the cumulative paid cost by date $i$, and let $b_i$ count labels that
already belong to groups of at least two labels.  A mixed scalar
retires its labels, which are retained as ghosts through $N$.

In the early range assume the actual propagated cover has at most
$CT_i$ paid pieces of diameter at most $Cs_i$, together with their
actual connections, where
\begin{equation}
 s_i\leq C\ell^{-(N-i)}(1+N-i)^{d_R},\qquad Cs_i<\tau/4.
 \label{ha:cover-scale}
\end{equation}
Each piece meets at most one copy.  A connected group meeting $h\geq2$
separated supports needs a spanning network of length at least
$c\tau(h-1)$.  Its paid cover constructs one of length at most
$Cs_iT_i$.  Therefore
\begin{equation}
 b_i\leq C\tau^{-1}s_iT_i.
 \label{ha:prefix-packing}
\end{equation}
This assertion is only early.  A single late cover can contain many
copies and need not satisfy it.

Suppose rough propagation after a first merger costs $e^{Ch}$ per date
for an $h$-label group, in the same original local/operator norm.
Choose $A_\tau=\lceil K_0(1+\log(1/\tau))\rceil$ and split at
$N-A_\tau$.  Then
\begin{equation}
 \sum_{i\leq N-A_\tau}b_i
 \leq\frac{C\ell^{-A_\tau}(1+A_\tau)^{d_R}}{\tau}T_N,
 \qquad
 \sum_{i>N-A_\tau}b_i\leq m(A_\tau+1).
 \label{ha:forest-sums}
\end{equation}
For sufficiently large fixed $K_0$, one reserved tree fraction pays
the entire early loss.  Hence
\begin{equation}
 e^{C\sum_i b_i}e^{-\kappa T_N}
 \leq C_0^m\tau^{-D_0m}e^{-\kappa T_N/2}.
 \label{ha:forest-result}
\end{equation}
No tree is re-spent at each ancestor merger.  The estimate applies
termwise to the original expansion with its original multiplicities;
it does not authorize summing a new family of labelled forests.

Only a small source ball may be available.  If a complete newborn map
is analytic for $\|X\|<R_j=cg_j^{\kappa_R}$ and bounded by $CR_j$,
its degree-$p$ ordered multilinear majorant costs
$C_1^pR_j^{1-p}\prod_l\|X_l\|$. Here the norm is the absolute
activity norm. In the $a_j^2$-normalized current-birth norm, replace
$R_j$ by $\widehat R_j=cg_j^{\kappa_R}/a_j^2$ throughout this bound.
The field-projection cost $a_j^{-2}$ cancels the $a_j^2$ norm of the
canonical reference. Thus unary propagation has no inverse-radius loss;
each new $p$-way merger loses $p-1$ radius powers, a fixed polynomial
coupling cost. Assign these powers
to distinct leaves: each child subtree retains one unassigned leaf;
at a merger assign all but one and pass the remaining representative
upward.  Every leaf is charged at most once.  If its first merger was
at $i_a\leq j$, then
$\log R_j^{-1}\leq C(1+N-i_a)$, whose sum is bounded by
$C\sum_i b_i+Cm$.  Equation~\eqref{ha:forest-result} pays it.
Existing Taylor/cluster symmetry factors remain those of the original
series; analyticity alone is not a claim of a factorial-free bound for
an arbitrary unordered coefficient.

A mixed scalar created at $l$ has root density $r_l^{-4}$.  Its ghosts
pay this, and an extra factor $e^{\lambda(N-l)}$ can be retained to
sum all primitive dates.  Distinct scalar factors in a square-free
term have disjoint labels.  A singleton deficit born before a merger
cannot be paid retroactively by that ghost: it needs the causal
singleton budget of Section~\ref{ha:scalars}.

\paragraph{Provenance.}
\path{work/merged_group_prefix_forest_2026-10-18.txt}
and the original rooted source majorant retained in
\path{work/initial_clone_census_quotient_audit_2026-10-18.txt}.

\section{From local chronological domination to normalized pressure}
\label{ha:pressure}

Let $F_a$ be complete actual scalar-composite copies on a finite torus,
and keep the original unmarked expectation $\mu_0$.  Restriction of the
copy family means setting its variables to zero in the same construction.
After exporting every universal marked scalar at birth, the exact identity is
\begin{equation}
 M(\boldsymbol\xi)=\mu_0\prod_a(1+\xi_aF_a)
     =e^{U(\boldsymbol\xi)}H(\boldsymbol\xi),\qquad
 H=\frac{Z_{\rm ren}(\boldsymbol\xi)}{Z_{\rm ren}(0)},
 \quad H(0)=1.
 \label{ha:pressure-identity}
\end{equation}
The denominator is the same source-free partition function.  The
history-dependent deficits in \eqref{ha:birth-scalar} are in $H$, not
in $U$.  If an original normalization was source dependent, its exact
numerator over the fixed reference inverse is retained before this
factorization.

\paragraph{Conditional local-to-global proposition.}
Assume the first-contact classification supplies the following summed
bound in the norm \eqref{ha:local-algebra}, with fixed weakened reserves:
\begin{equation}
 \sum_S t_S\|T^{\rm ren}_{H,S}f_{\rm test}\|
       \leq B_H^{\#}\|f_{\rm test}\|
 \label{ha:spectator-domination}
\end{equation}
for every original analytic spectator, together with the corresponding
rooted bulk and boundary bounds.  The empty coefficient is unchanged.
Suppose the same positive family retains each copy-to-output carrier.
Then, for terminal volume $V_T$ and $m$ copies,
\begin{equation}
 \|H\|_{\mathbf t}\leq e^{C_pV_T},
 \qquad
 \|U-\sum_a\xi_au_a\|_{\mathbf t}\leq C_vm,
 \quad u_a=[\xi_a]U.
 \label{ha:PV}
\end{equation}

\emph{Proof of the first bound.}
Linear integrations and restrictions act coefficientwise. Multiplication
is submultiplicative in the finite label algebra; exponentials and Mayer
differences are bounded by their original positive series; chronological
composition uses \eqref{ha:spectator-domination}.  The original final
stability proof therefore yields
$\|Z_{\rm ren}(\boldsymbol\xi)\|_{\mathbf t}\leq e^{C_+V_T}$.
Its source-free lower bound is
$Z_{\rm ren}(0)\geq e^{-C_-V_T}$ in the same normalization.  Division
gives $C_p=C_++C_-$.  No separate linked-cluster theorem for a normalized
live conditional law is needed.  Conversely, a family estimate only
on $T_H1$ or only on nonvacuum coefficients would not justify this step.

\emph{Proof of the mixed-scalar bound.}
Anchor a mixed universal birth to one original marked vertex of a
specified copy $a$, in a fixed compact support $K$, and retain its
root-to-output path.  At current birth date $l$ the output-root count
costs at most
\begin{equation}
 C_Kr_l^{-4}P(\gamma_l)(1+T)^{d_0}
 \label{ha:scalar-root-count}
\end{equation}
times its original rooted carrier majorant.  This is a current root
count, not another sum over old fusion roots.  The polynomial in the
carrier is paid by a spare exponent.  In the prefix estimate include
the ghost charge
$\exp[(4\log\ell+C_P+\lambda)(N-l)]$, where
$P(\gamma_l)\leq Ce^{C_P(N-l)}$ and $\lambda>0$ is fixed.
It pays \eqref{ha:scalar-root-count} and leaves
$Ce^{-\lambda(N-l)}$.  All primitive dates are now summable.

The remaining original rooted series sums unmarked connectors and
relative marked roots with the fixed reduced reserve.  Separated copy
positions have density at most a fixed power of $\tau^{-1}$ in a
carrier enlargement; retain the corresponding carrier moment.  Choose
the final label weight with an additional fixed factor $\tau^D$ to
absorb this density, profile derivatives, and the forest loss.
Its local sum is a small constant $c$, independent of the number of
copies.  Selecting an anchor in a coefficient with label set $S$ costs
at most $|S|\leq2^{|S|}$, paid by reducing $c$ once while retaining
$c^{|S|}$.  Thus
\begin{equation}
 \sup_a\sum_{S\ni a,\ |S|\geq2}t_S|U_S|\leq C_v.
 \label{ha:rooted-vacuum}
\end{equation}
Summing over $a$ proves the second bound of \eqref{ha:PV}.  This proof
uses the \emph{same retained rooted family} as
\eqref{ha:spectator-domination}; separate bounds on each fixed subset
would not suffice.  Universal singletons are excluded and may be
extensive. \hfill$\square$

An explicit sufficient version of \eqref{ha:spectator-domination} is a
coefficientwise comparison, after complete-copy cancellations and causal
scalar extraction, with the original simultaneous source family at
radii $r_a=P(G)^{-1}$:
\begin{equation}
 t_S\operatorname{Maj}(\hbox{decorated term}_S)
 \leq c^{|S|}r_S\operatorname{Maj}_0
                       (\hbox{same term; weakened reserves}),
 \qquad t_a\varepsilon(G)\leq cr_a.
 \label{ha:license}
\end{equation}
Original scalar derivatives are extracted at the simultaneous local
polydisc, with at most five internal labels per copy.  One must not
differentiate a history-diameter-dependent Cauchy radius once for each
label.  The first-contact lanes and forest licenses aim to prove
\eqref{ha:license}; the remaining source-admission gap prevents treating
this application as an unconditional theorem.

\subsection{Compatible centering and actual arrays}

The singleton residual mean is controlled without a local partition
ratio.  Choose a fixed terminal-volume cell of volume $v_0$, compatible
with all relevant past partitions and large enough for a buffered copy.
Disseminate $D=F-u$ over its $B$ reflected cells.  Restriction consistency
and compatible covariance identify the singleton scalars as the
disseminates of $u$.  By \eqref{ha:PV}, the top coefficient is bounded by
$t^{-B}\exp(C_pBv_0+C_vB)$.  The chessboard inequality and its $B$th
root give
\begin{equation}
 |\mu_0F-u|\leq C_0t^{-1},\qquad C_0=e^{C_pv_0+C_v}.
 \label{ha:residual-mean}
\end{equation}
There is no zero-free logarithm of $H$ in this argument.

Put $F_a^c=F_a-\mu_0F_a$.  Nilpotence gives exactly
\begin{equation}
 \mu_0\prod_a(1+\xi_aF_a^c)
 =\exp\!\left(U_{\rm mix}+\sum_a\xi_a(u_a-\mu_0F_a)\right)H.
 \label{ha:exact-centering}
\end{equation}
Equations~\eqref{ha:PV}--\eqref{ha:residual-mean} bound its full coefficient
norm by $\exp[C_pV_T+(C_v+C_0)m]$.  Top coefficient extraction costs
exactly $\prod_at_a^{-1}$, with no extra factorial.

Place a time array of $m$ copies in one rectangular cell whose spatial
cross-section depends only on their fixed supports.  Its terminal volume
is at most $C_Km$.  Disseminate the \emph{whole array cell} over a large
torus, apply the preceding bound, and take the root by the number of
cells.  The resulting chessboard seminorm is
\begin{equation}
 \left\|\prod_{a=1}^mF_a^c\right\|_{\rm CB}
                       \leq\prod_{a=1}^m(C_Kt_a^{-1}).
 \label{ha:array-bound}
\end{equation}
Fixed-regulator polynomial endpoint vectors are expanded into link-entry
monomials and assigned to other closed cells before and after the array.
Their sup norms give a prefactor independent of $m$.  Enlarging the
array's transverse cross-section to contain arbitrary endpoints would
instead produce an inadmissible endpoint-dependent constant per copy.
Reflection doubling removes the fixed prefactor in the thermodynamic
state.  Complex insertions stay intact; taking their pointwise absolute
values or replacing a complex source by its real part is not part of
the proof.

\section{Causal prefix periods and the conditional copied estimate}
\label{ha:copies}

The centering identity does not make arbitrary translates of a
prefix-subtracted fusion into compensated regular prototypes.  A
translation can shift its numerical lower-subset coefficient pattern
relative to the source scaffold.  Full centering corrects only its
identity coefficient.

Use instead the common causal prefix period of the two adjacent frames,
including their scalar convention and all partition choices through
fusion.  Its physical size obeys
\begin{equation}
 p_G\leq Cr_f(1+\log G)^d\leq e^{-cG}P(G).
 \label{ha:prefix-period}
\end{equation}
The terminal matching matrices are spatially constant and do not enlarge
this period.  A period translate is exactly the same unweighted prototype
with translated envelope.  Future kernels need not commute with that
translation: complete hat grouping and its exact profile commutator are
performed on the fixed future trajectory before norms.

The common origin reflection preserves the actual old-$R$ regular base,
retained coordinates, scalar-radius windows, gauge/coarea factors, and
intermediate normalizers.  Its finite products intertwine with the
future-corrected projection, so a zero corrected coordinate remains zero.
It is sufficient to adjoin finitely many actual reflected prototypes;
the selected twenty-one columns need not already form a reflection
representation.  A plane on an even $p_G$ grid is an origin reflection
followed by a permitted translation.

The residual-mean cell used in \eqref{ha:residual-mean} still uses the
larger all-past-compatible period.  The long array uses the finer
even-$p_G$ grid.  For
$G\geq A[1+\log_+(1/\tau)]$, rounding its faces within empty buffers
costs $O(p_G)\ll\tau$ and keeps volume $O(m)$.  At fixed bare regulator
the free torus size factor can be chosen divisible by any fixed finite
array dimensions and the source period, with even counts.  Take volume
limits before arbitrary doubling and then the ultraviolet subsequence.

Combining the preceding conditional estimates gives the actual conclusion
of this part of the argument.  For a fixed compact center support and a
fixed complete adjacent-frame difference $E_G$, there are finite
$B,D,S$ such that
\begin{equation}
 \|\mathcal D_\tau(E_G(f)^c)\|
 \leq C\tau^{-D}G^B
       \{e^{-c(\log G)^s}+G^{-K}\}\,p_S(f),\qquad s>1.
 \label{ha:copied-estimate}
\end{equation}
Here $\mathcal D_\tau$ denotes the bounded slab insertion with positive
padding, not an already defined undamped operator.  The allowed array
copies are the causal-period translates and finite reflections just
specified.  Constants are independent of bare depth and array length.

The parameter order is finite and noncircular.  First fix the admission
exponent $\kappa_R$ above the original stability threshold, for example
at least eight, and all geometric and analytic reserves.  Fix the fusion
order, creation and matching powers, rough propagation constants, ghost
charge, and finite test seminorm.  These determine $B,D$.  Then choose
the source-free native output power so that the remaining $K$ exceeds
$B$ and any later finite domain/locality losses.  Increase the weak
floor last.  Native $R''$ remains a young renewal and receives no false
fresh bad-field factor.  The weak terminal band is held fixed.  The
proof does not provide one source radius uniform over arbitrary terminal
inverse coupling tending to infinity.

\paragraph{Provenance.}
\path{work/copied_fusion_pressure_consumer_2026-10-18.txt},
\path{work/prefix_period_copied_slab_consumer_2026-10-18.txt},
and \path{work/copy_pressure_normalization_integration_audit_2026-10-18.txt}.
The arbitrary-bare-translate wording in the older pressure consumer is
superseded by the restricted causal-period statement above.

\section{Two global colours and the corrected Holder packing power}
\label{ha:two-colour}

The following is a different, degree-two problem.  The copied bound
\eqref{ha:copied-estimate} cannot be used as though it supplied a rough
two-colour radius $R^\theta/P(G)$; its padding power may be much larger.
Two global signed unions of packets are labelled by
$z_1,z_2$, $z_1^2=z_2^2=0$, only after each complete internal scalar
fusion coefficient has been formed.  A prefusion operator has the actual
coefficients
$T(W_1\,\cdot)$, $T(W_2\,\cdot)$, and $T(W_1W_2\,\cdot)$.
The last is not the product of its values on the two separate marks.

For an ordinary vertex $V=V_0+z_1V_1+z_2V_2+z_1z_2V_{12}$,
\begin{equation}
 [e^V-1]_a=e^{V_0}V_a,\qquad
 [e^V-1]_{12}=e^{V_0}(V_{12}+V_1V_2).
 \label{ha:vertex-two}
\end{equation}
For the normalized regular conditional map $K$ the logarithmic mixed
coefficient is
$KV_{12}+K(V_1V_2)-(KV_1)(KV_2)$.  These are exact identities with
the original denominators.  The original source connected series, with
its original symmetry factors and chronological order, gives
\begin{equation}
 [R'_j(X)]_{12}=\sum_{\mathcal C}\varphi^T(\mathcal C)
 \left\{\sum_vp_{v,12}\prod_{u\ne v}p_{u,0}
  +\sum_{v\ne w}p_{v,1}p_{w,2}\prod_{u\ne v,w}p_{u,0}\right\}.
 \label{ha:two-cluster}
\end{equation}
For operator factors the products and replacements have their original
order.  No division by $p_{v,0}$ is used.  A quadratic marked-position
count is paid by a fixed reserve.

The required input is a \emph{projective profile representation} before
relative positions are summed:
\begin{equation}
 p_{v,a}[F]=\sum_\alpha L_{v,a,\alpha}\lambda_{v,\alpha}(F),
 \quad
 p_{v,12}[F_1,F_2]=\sum_{\alpha,\beta}
  L_{v,12,\alpha\beta}\lambda_{v,\alpha}(F_1)
                           \lambda_{v,\beta}(F_2).
 \label{ha:slots}
\end{equation}
The slots are evaluations, normalized hats, and original endpoint
differences with their carrier moments.  The sums of coefficient
\emph{operator} norms, with relative positions inside the current-rooted
majorant, must be bounded.  A supremum estimate for each profile or a
homogeneous matrix estimate does not imply this norm.

Suppose an actual opened history has the first- and mixed-input bounds
$A_a[J_H+\sum_nP(\gamma_n)N_n]$ and
$A_1A_2[J_H+\sum_nP(\gamma_n)N_n]^2$ in that norm.  Here the $N_n$
are the exact paid own-layer census, not a newly inserted number of old
cells, and $J_H$ is a current carrier moment.  The retained occupation
reserve gives
\begin{equation}
 \left(1+\sum_nP(\gamma_n)N_n\right)^2
               e^{-c\sum_nq_nN_n}\leq C.
 \label{ha:degree-two-absorb}
\end{equation}
For example the original stretched-log $q_n$ dominates
$\log P(\gamma_n)$ with a fixed margin; expanding the degree-two
polynomial and assigning reserve to its at most two chosen factors
proves the bound.  The two colours may use the same cell.  Apply this
inequality before the existing rooted history sum.  Equations
\eqref{ha:two-cluster}--\eqref{ha:degree-two-absorb} then bound a new mixed
birth by $C(B_{12}+B_1B_2)$ with one fixed weaker reserve.  A source-free
native variation supplies one changed native vertex, not a native factor
per colour.

Causal scalar removal is still exact.  If
$V=z_1V_1+z_2V_2+z_1z_2V_{12}$, multiplying the density by $e^{-V}$ sends
\begin{equation}
 \rho_{12}\longmapsto\rho_{12}-V_1\rho_2-V_2\rho_1
                            +(V_1V_2-V_{12})\rho_0.
 \label{ha:two-scalar}
\end{equation}
An existing mixed coefficient is carried as such; it is not reset to a
product of first coefficients.  Its current-date norm and the prefusion
seed reserve are substantive inputs.  Freezing the maps and centering
numbers proves the algebra but not that quantitative assertion.

For separated global supports at distance at least $cR$, a connected
mixed birth at $u$ retains a path meeting both supports, including paid
component interiors.  Its tree cost satisfies
$T_u\geq cR/(r_uMR_u)$.  If
$r_w=R\gamma_w^{-A}$ and $u=w-t$, bounded drift gives
\begin{equation}
 \frac{R}{r_uMR_u}\geq
 c\frac{\gamma_w^A}{(1+\log\gamma_w)^d}\ell^{t/2}.
 \label{ha:early-path}
\end{equation}
Every fixed polynomial in the early coupling and age is summable against
the resulting exponential, uniformly in the initial fusion $G$.
Universal mixed cohorts retain fixed relative offsets and their finite
phase period.  For a nonnegative offset majorant with $p^4$ phases,
$\sum_\delta\max_\sigma a_{\delta\sigma}
\leq p^4\max_\sigma\sum_\delta a_{\delta\sigma}$.
The period conversion is a current polynomial.  It does not turn a live
operator into an activity or prove its spectator norm by evaluation at $1$.

\subsection{One original envelope per linear lineage}

Fix $0<\theta<1/2$.  For the original smooth profile define
$A_dF(x)=(F(x+d)-F(x))/|d|^\theta$.  The elementary inequality
\begin{equation}
 |\Delta_h\Delta_dF(x)|\leq
 2[F]_{2\theta}\min(|h|^{2\theta},|d|^{2\theta})
 \leq2[F]_{2\theta}|h|^\theta|d|^\theta
 \label{ha:holder-double}
\end{equation}
implies $\|A_dF\|_{C^\theta}\leq C\|F\|_{C^{2\theta}}$ uniformly in
$d\ne0$, with the usual sup estimate for $|d|\geq1$.  There is a cost
$R^{-2\theta}$ per colour for a packet of size $R$, not the earlier
unjustified $R^{-\theta}$ cost.

This regularity is paid only once.  If $L_j$ is the frozen linear map on
the complete dated operator/scalar ledger and $M_{F_j}$ multiplication
at its current root, the exact commutator identity is
\begin{equation}
 E_{j+1}=L_jE_j+
       (L_jM_{F_j}-M_{F_{j+1}}L_j)L_{j:i},
 \qquad E_j=L_{j:i}M_{F_i}-M_{F_j}L_{j:i}.
 \label{ha:fixed-lineage}
\end{equation}
An existing gradient lineage is transported linearly, not refrozen as a
new envelope at each date.  The new forcing contains one original
profile difference.  At the first bilinear merger the product of the
two slots is a fixed envelope for its mixed lineage.  Later operations
linear in that mixed coefficient keep it and its origin tag.  With
only two colours there is no unbounded sequence of normalized profile
derivatives.

In physical units a difference born at $l$ and read at $j\geq l$ costs
$C[F]_{2\theta}r_j^\theta r_l^\theta
\leq C[F]_{2\theta}r_j^{2\theta}$, with its original carrier moment.
Thus the conditional current mixed birth bound has profile cost
\begin{equation}
 C P(\gamma_j)(1+R_j^{-2\theta})^2,\qquad R_j=R/r_j.
 \label{ha:profile-cost}
\end{equation}
Packing $O(R_N^{-4})$ pairs and dividing the current root count by that
number supplies $R_j^4$.  For $R_j\leq1$ the corrected late and terminal
power is
\begin{equation}
             P(\gamma_j)R_j^{4-4\theta}.
 \label{ha:packing-correct}
\end{equation}
For instance $\theta=1/4$ gives power three.  The original envelope,
causal scalar tags, current output-root convention, and retained relative
paths are all needed for this calculation.

This is a conditional current-birth attachment.  It does not yet sew the
early ten-scalar Gaussian comparison, late universal sum, and normalized
terminal remainder into a full-state relative Gram lower bound.  In
particular, \eqref{ha:packing-correct} is not a statement that the
twenty-one limiting fields are independent.

\paragraph{Provenance.}
\path{work/two_colour_current_birth_kernel_2026-10-18.txt},
especially its exact input in Sections 6--7 and its corrected
fixed-lineage construction in Sections 10--11; see also
\path{work/current_date_two_mark_gram_budget_2026-10-18.txt}.

\section{Matching actual scalar frames without a Gram theorem}
\label{ha:matching}

Let $V=\operatorname{Sym}^2(\Lambda^2\mathbb R^4)$, of dimension twenty-one.
Choose a finite set of scalar packet configurations, of orders at most
five, whose normalized open Gaussian coefficient matrix $A$ is invertible.
This is local coefficient rank only.  For each column retain the complete
density germ
\begin{equation}
 D_{[k]}\left[e^{-V_{\rm prefix}(z)}\mathcal O_{\rm source}(z)\right]_{z=0}.
 \label{ha:density-frame}
\end{equation}
Its prefix scalar convention is causal and frozen at its fusion date.
If the conditional logarithm has coefficients $W_B$, the density
coefficient is exactly $\sum_{\mathcal P}\prod_{B\in\mathcal P}W_B$.
After the order-$k$ normalization, a partition with $t$ blocks costs
\begin{equation}
 C_k b_f^2\left(\frac{R_f^4b_f^2}{g_s^2}\right)^{t-1}.
 \label{ha:density-partition}
\end{equation}
The bracket tends to zero at the chosen packet scale.  Products with
$t\geq2$ have field degree at least four and no quadratic marginal.
Thus the actual density frame can satisfy both
$\Pi W=A+o(1)$ and $\|W\|_b/b^2\leq C$.  A merely polynomial original
norm would not give the latter statement.

Assume the physical regular graph has coordinates $(c,r)$, graph $T_f$,
and future-corrected projection
$\pi_f^N(c,r)=c+L_f^N(r-T_fc)$ with
$\|T_f\|\leq Cg_f^2$, $\|L_f^N\|\leq C\eta_f$.
The bounded physical frame gives
\begin{equation}
 \|C_f^N-A\|\leq C(o(1)+\eta_f),\qquad C_f^N=\pi_f^NW_f.
 \label{ha:matrix-inverse}
\end{equation}
For a sufficiently weak floor it is invertible with uniformly bounded
inverse.  No OS Gram lower bound has been used.

For actual frames born at $i\leq f$, let $P_{i,f}$ be the marginal
transport and set
$B_i=P_{i,N}C_i^N$, $B_f=P_{f,N}C_f^N$.  Regard the frames as row
vectors of actual observables.  The difference
\begin{equation}
 E_{i,f}=F_iB_i^{-1}-F_fB_f^{-1}
 \label{ha:actual-difference}
\end{equation}
has zero corrected coordinate at the common date $f$, because
\begin{equation}
 \pi_f^N(K_{i,f}W_iB_i^{-1}-W_fB_f^{-1})
            =P_{f,N}^{-1}-P_{f,N}^{-1}=0.
 \label{ha:matching-zero}
\end{equation}
Matrix order matters; no commutativity is assumed.  Products and inverses
have their fixed polynomial coupling bounds.  Every column remains a
linear combination of actual complete scalar germs with their original
chronological numerators, scalar deficits, and reference inverses.
No abstract graph representative is inserted into a mixed history.

The complete unweighted center cohort is grouped and matured before
projection, using the exact tent residue balance.  This yields a constant
matching matrix.  For a nonconstant center profile the difference from
its current anchor is an explicit source-gradient family.  Period averaging
does not by itself eliminate this family or give arbitrary-translate
compensation; Section~\ref{ha:copies} states the narrower copied class.

\paragraph{Provenance.}
\path{work/actual_scalar_fusion_frame_matching_2026-10-18.txt}.

\section{A conditional heat-damped limit on one scalar Hilbert space}
\label{ha:heat-limit}

Assume \eqref{ha:copied-estimate} for adjacent normalized actual frames,
and the scalar Hilbert-space inputs of Section~\ref{ha:scalar}.  Let
$G_n=2^nG_0$ and let $H_{G_n}^c(f)$ denote the centered complete frame.
For $\operatorname{supp}f\subset\{|t|<T\}$, write
\begin{equation}
 B_{G,a}(f)=e^{-aH}H_G^c(f)e^{-aH},\qquad a>T,
 \label{ha:finite-heat}
\end{equation}
as notation for the bounded ordered slab form.  This does not assume an
undamped operator exists.  The difference estimate is
\begin{equation}
 \|B_{G,a}(f)-B_{2G,a}(f)\|
 \leq C\min(1,a-T)^{-D}d(G)p_S(f),
 \quad d(G)=G^B\{e^{-c(\log G)^s}+G^{-K}\}.
 \label{ha:dyadic-difference}
\end{equation}
The finite-scale scalar anchor costs only
\begin{equation}
 \|B_{G,T+\tau}(f)\|\leq Ce^{CG}\min(1,\tau)^{-D_0}p_{S_0}(f).
 \label{ha:heat-anchor}
\end{equation}
There are at most five scalar factors, with fixed distinct internal
time coordinates and positive gaps whose inverses cost $e^{CG}$.
All deterministic coefficients and finitely many derivatives have the
same type of cost.  Spatial separation alone would not justify this
ordered heat product.  The finite configurations can be perturbed to
distinct times while preserving a positive singular-value margin.

At each fixed finite list of $G_n$, take volume limits first, pass the
bounded slab forms to the scalar continuum Hilbert space by their ordered
matrix elements, and fix the finitely many allowed date offsets on a
subsequence.  Matching matrices are bounded at fixed $G_n$; smooth
coefficient functions have bounded fixed derivatives on compact centers.
A countable diagonal subsequence therefore puts all scales on one
Hilbert space and passes adjacent matching identities together.  There
is no assertion of operator-norm convergence between different bare
Hilbert spaces, nor of uniqueness of these limiting matching data.

The dyadic tail satisfies, for a slightly smaller $c'$,
\begin{equation}
 \sum_{k\geq n}d(G_k)
 \leq C\{e^{-c'(\log G_n)^s}+G_n^{B-K}\}=:d_*(G_n).
 \label{ha:dyadic-tail}
\end{equation}
The power term is geometric; superlinear decay in $\log G_k$ absorbs
its polynomial factor and the remaining sum.  Thus $B_{G_n,a}$ converges
in the norm of maps from one fixed finite test seminorm to bounded
operators.  Denote the limit by $B_a(f)$.  It is centered and satisfies
\begin{equation}
 B_b(f)=e^{-(b-a)H}B_a(f)e^{-(b-a)H},\qquad b>a>T.
 \label{ha:heat-compatible}
\end{equation}

For small $\tau$, take the first dyadic $G$ exceeding
$A[1+\log(1/\tau)]$.  Then $e^{CG}$ in
\eqref{ha:heat-anchor} is only a fixed power of $\tau^{-1}$, and
\eqref{ha:dyadic-tail} bounds the rest.  Consequently
\begin{equation}
 \|B_{T+\tau}(f)\|\leq C\tau^{-D'}p_S(f).
 \label{ha:heat-bound}
\end{equation}
For tests supported in $[-\tau,\tau]\times K$, choose the same anchor
so that the actual stencil radius and prefix period are below $\tau/10$.
The same proof, on one fixed containing compact, gives the uniform
small-slab version
\begin{equation}
 \|B_{4\tau}(f)\|\leq C_K\tau^{-D'}p_S(f).
 \label{ha:small-slab}
\end{equation}
A fixed positive support collar cannot simply be set to zero; this
anchor choice is the needed justification.

On $\mathcal D_{\rm fin}=\bigcup_E1_{[0,E]}(H)\mathcal H$ define
\begin{equation}
 q_f(u,v)=\langle e^{aH}u,B_a(f)e^{aH}v\rangle.
 \label{ha:heat-form}
\end{equation}
It is independent of $a$ by \eqref{ha:heat-compatible}.  Its spatial and
Euclidean-time translation identities follow from the exact causal-period
identities, since $p_{G_n}\to0$.  Spatially, heat slack and the scalar
cone imply $\|(U(y)-1)e^{-\delta H}\|\leq C|y|/\delta$.
For time shifts, use the bounded unequal-padding expression
\begin{equation}
 e^{-(a-a_0-t)H}B_{a_0}(f)e^{-(a-a_0+t)H}
 \label{ha:unequal-padding}
\end{equation}
with both times positive, and
$\|He^{-\delta H}\|\leq C/\delta$.  Approximate a fixed shift by prefix
period shifts and pass the finite-seminorm bounds.  Thus
$q(f_t)=e^{tH}q(f)e^{-tH}$ on finite-energy vectors.  The argument is
translation-specific and supplies no continuous rotation covariance.

\paragraph{Provenance.}
\path{work/dyadic_actual_frame_heat_limit_2026-10-18.txt}.

\section{From heat forms to closable tempered translation fields}
\label{ha:operator}

Polynomial heat-sandwich bounds alone do not give an undamped operator.
For example, with $He_n=ne_n$ on $\ell^2$, the formal matrix $q_{mn}=1$
has a bounded rank-one heat sandwich of norm $O(\tau^{-1})$, but its
application to a basis vector is not square summable.  The exact
Euclidean-time identity obtained above is the extra input used here.

\subsection{A finite-band distribution ODE and its uniform inverse}

Fix a spatial test $g$ and $P_E=1_{[0,E]}(H)$.  The compressed form
$P_Eq(\chi\otimes g)P_E$ is a distribution in time with values in
$\mathcal B(P_E\mathcal H)$.  The generator $H_E$ is bounded.  Its
translation law implies the distribution equation
$\partial_tT=\operatorname{ad}(H_E)T$.  Multiplying by
$e^{-t\operatorname{ad}(H_E)}$ gives a distribution of zero derivative,
which is a constant times Lebesgue measure: every compact test with
integral zero is the derivative of a compact smooth test.  There is
therefore a unique bounded operator $A_E(g)$ such that
\begin{equation}
 P_Eq(\chi\otimes g)P_E
    =\int\chi(t)e^{tH_E}A_E(g)e^{-tH_E}\,dt.
 \label{ha:distribution-ode}
\end{equation}
Uniqueness gives $P_{E'}A_E(g)P_{E'}=A_{E'}(g)$ for $E'<E$.
This obtains a time-zero value after compression, not by evaluating
an arbitrary distribution at a point.

Choose one even nonnegative $\chi\in C_c^\infty((-1,1))$ with integral
one, put $\chi_\tau(t)=\tau^{-1}\chi(t/\tau)$, and define
$\Lambda(s)=\int\chi(t)e^{ts}\,dt$.  It is positive on the real line.
For $\tau=(1+E)^{-1}$, all spectral differences $\tau(E_1-E_2)$ lie
in $[-1,1]$.  Extend $1/\Lambda$ from this interval to a fixed smooth
compactly supported function $m$ on $\mathbb R$.  Its Fourier transform
has finite $L^1$ norm.  Consequently
\begin{equation}
 \mathcal M_{E,\tau}(B)=\int\widehat m(u)
                e^{iu\tau H_E}Be^{-iu\tau H_E}\,du
 \label{ha:laplace-inverse}
\end{equation}
is bounded uniformly in $E$ and in the dimension of $P_E\mathcal H$,
and inverts $\Lambda(\tau\operatorname{ad}H_E)$.
One can verify the inverse identity first on finite spectral step
approximations to bounded $H_E$, then pass in operator norm using the
$L^1$ Fourier majorant.  No finite-rank energy projection is assumed.

The small-slab estimate \eqref{ha:small-slab} and
$p_S(\chi_\tau\otimes g)\leq C\tau^{-(S+1)}p_S(g)$ imply
\begin{equation}
 \|A_E(g)\|\leq C_K(1+E)^p p_S(g)
 \label{ha:band-polynomial}
\end{equation}
for a fixed integer $p\geq D'+S+1$.  Undoing $4\tau$ of heat on each
side costs at most $e^{8\tau E}\leq e^8$ before applying
\eqref{ha:laplace-inverse}.  Inverting a fixed positive time average
instead would lose an exponential in $E$.

Spatial translation covariance makes the constant uniform on translates
of a fixed unit compact.  A smooth unit-cube partition and the weight
$(1+|x|)^4$, which is summable in spatial dimension three, extend
\eqref{ha:band-polynomial} to spatial Schwartz tests with
\begin{equation}
 p_{S,4}(g)=\sum_{|\alpha|\leq S}
              \sup_x(1+|x|)^4|\partial^\alpha g(x)|.
 \label{ha:spatial-seminorm}
\end{equation}
The compatible $A_E$ define a form $A(g)$ on finite-energy vectors.
It need not be an operator there.  Unit-band Schur estimates do give
$\|e^{-\tau H}A(g)e^{-\tau H}\|
\leq C\tau^{-(p+1)}p_{S,4}(g)$.
Equation~\eqref{ha:distribution-ode}, with positive left and right
paddings, identifies its Euclidean one-insertion integral with the
original $B_a(f)$ by equality of all finite-energy compressions.

\subsection{Real-time smearing and a common invariant domain}

Let $P_n=1_{[n,n+1)}(H)$ and write
$A_{mn}(g)=P_mA_E(g)P_n$ for sufficiently large $E$.  Consistency makes
this independent of $E$, and
$\|A_{mn}(g)\|\leq C(1+m+n)^pp_{S,4}(g)$.
For a spacetime Schwartz test define
\begin{equation}
 \Phi_{mn}(f)=\int_{\mathbb R}
                e^{itH_m}A_{mn}(f(t,\cdot))e^{-itH_n}\,dt.
 \label{ha:real-smearing}
\end{equation}
The integral converges in norm.  Write $H_m=m+K_m$,
$H_n=n+K_n$, with $\|K_m\|,\|K_n\|\leq1$.  Integration by parts in
$e^{it(m-n)}$, with derivatives falling on bounded residual generators
and the test, proves for any fixed $k$
\begin{equation}
 \|\Phi_{mn}(f)\|\leq C_k(1+m+n)^p
                 (1+|m-n|)^{-k}\mathcal N_k(f),
 \label{ha:block-decay}
\end{equation}
where one may take
\begin{equation}
 \mathcal N_k(f)=\sum_{j\leq k,|\alpha|\leq S}
 \int dt\,\sup_x(1+|x|)^4
                   |\partial_t^j\partial_x^\alpha f(t,x)|.
 \label{ha:block-seminorm}
\end{equation}
A time weight of order two bounds this by an ordinary finite Schwartz
seminorm.  Compact Fourier support is unnecessary.

Since $1+m+n\leq C(1+n)(1+|m-n|)$, block Schur summation for
$k>p+1$ yields
\begin{equation}
 \|\Phi(f)(1+H)^{-p}\|\leq C\mathcal N_k(f).
 \label{ha:one-sided-energy}
\end{equation}
More generally, for every integer $r\geq0$, choose $k>p+r+1$ to get
\begin{equation}
 \|(1+H)^r\Phi(f)(1+H)^{-(p+r)}\|
                  \leq C_r\mathcal N_{k}(f).
 \label{ha:all-weight-energy}
\end{equation}
The consistent blocks define $\Phi(f)$ on $\operatorname{Dom}(1+H)^p$,
and in particular on the dense invariant space
$\mathcal D_\infty=\bigcap_r\operatorname{Dom}(1+H)^r$.
Equation~\eqref{ha:all-weight-energy} proves invariance and continuous
Schwartz dependence on this space.  Temporal and spatial translations
are exact.  Their positive-energy representation is the already supplied
scalar one; no Lorentz transformation law of the new label is added.

To prove closability, adjoin the actual reflected/conjugate finite frame
before taking limits.  It is enough to enlarge the family finitely; one
need not identify it with a prescribed twenty-one-component tensor
reflection.  The even real averaging test above and uniqueness in
\eqref{ha:distribution-ode} give
$A_E(v,g)^*=A_E^\dagger(\bar v,\bar g)$.  Thus
\begin{equation}
 \Phi_v(f)^*\supset\Phi_{\bar v}^{\dagger}(\bar f)
                   \quad\hbox{on }\mathcal D_\infty.
 \label{ha:adjoint}
\end{equation}
The adjoint has a dense domain, so each smeared field is closable.
Hermitian combinations with its companion are symmetric on the common
domain.  These conclusions concern genuine real-time smeared operators,
not merely bounded heat forms.

\paragraph{Provenance.}
\path{work/heat_form_to_operator_distribution_2026-10-18.txt}.
The original scalar reconstruction's polynomial energy estimate and the
new frame's band inversion are different arguments, with their own
identification hypotheses explicitly retained.

\section{Locality from the shrinking scalar stencils}
\label{ha:locality}

The field-domain theorem does not by itself imply locality.  Here the
actual finite scalar-stencil provenance supplies a controlled crossing
argument.  Assume the scalar field and its adjoint have the all-weight
energy bounds and locality on $\mathcal D_\infty$.  Assume in addition
that the identified, spatially smeared scalar time-zero form $a(g)$,
supplied by the full-point ordered-insertion construction, satisfies
the following bound for one fixed $q$:
\begin{equation}
 \|(1+H)^{-q}a(g)(1+H)^{-q}\|\leq Cp_s(g).
 \label{ha:scalar-form-bound}
\end{equation}
Assume also the copied convergence already used above, and that the
complete stencil diameter and prefix period obey
$\delta_G+p_G\leq e^{-cG}P(G)$, while inverse internal time gaps and
finitely many coefficient/test derivatives cost at most $e^{CG}$.

\subsection{A downstream averaging and quantitative approximation}

Average the actual finite-$G$ form over its prefix torus, translating
back by the Hilbert-space translations.  The normalized average removes
the periodic center dependence of the deterministic coefficients and
leaves a continuously translated scalar-stencil expression.  It is
performed \emph{after} the copied estimate, with positive heat slack in
Euclidean time.  It makes no assertion that arbitrary translates were
admissible producer clones.  Since the limit is translation covariant,
this probability average has the same error $d(G)$ and the same limit.

Its time-zero form is a finite sum or nuclear integral of ordered terms
\begin{equation}
 e^{s_1H}a(g_1)e^{-(s_1-s_2)H}a(g_2)\cdots
                           a(g_k)e^{-s_kH},
 \quad k\leq5,\quad s_1>\cdots>s_k,\quad |s_a|\leq\delta_G.
 \label{ha:stencil}
\end{equation}
All internal heat gaps stay positive.  The absolute coefficient and
fixed seminorm integral costs at most $e^{CG}$.  Joint spatial kernels
need not have finite tensor rank: rescale their finitely many packet
boxes, expand in a product Fourier series, and integrate by parts a
fixed number of times in each spatial slot.  The resulting nuclear
sum is absolutely convergent in the required seminorm, with the same
type of $e^{CG}$ cost.

The small-time inverse applied to the difference gives
\begin{equation}
 \|P_E(A_G^{\rm av}(g)-A(g))P_E\|
             \leq C(1+E)^p d(G)p_s(g)
 \label{ha:band-error}
\end{equation}
when $G\geq C[1+\log(1+E)]$ and the stencil is much smaller than
$(1+E)^{-1}$.  There is also the crude bound
$\|e^{-\rho_GH}A_G^{\rm av}(g)e^{-\rho_GH}\|
\leq e^{CG}p_s(g)$ for $\rho_G\leq e^{-cG}P(G)$.
Choose $a_G=c_0G^{-1/2}$ and $E_G=L_0G/a_G$, with $L_0$ fixed large.
Band Schur summation bounds the high-energy part by
$Ca_G^{-1}e^{CG-ca_GE_G}$; the limiting form has an analogous polynomial
tail.  The low-energy part uses \eqref{ha:band-error}.  Hence for a fixed
$D_1$,
\begin{equation}
 \|e^{-a_GH}(A_G^{\rm av}-A)e^{-a_GH}\|
                   \leq C\{G^{D_1}d(G)+e^{-c'G}\}p_s(g).
 \label{ha:shrinking-heat-error}
\end{equation}
The same holds after real-time center smearing.  The small error has
not been multiplied by $e^{a_GE_G}$ or by $e^{CG}$.

\subsection{Gaussian time averaging controls the crossing}

Let $B$ and its adjoint partner satisfy the all-weight energy bounds,
and put $\alpha_t(B)=e^{itH}Be^{-itH}$.  For the normalized Gaussian
$k_\epsilon(t)$ of variance $\epsilon^2$, set
\begin{equation}
 B_\epsilon=\int k_\epsilon(t)\alpha_t(B)\,dt,\qquad
 \alpha_z(B_\epsilon)=\int k_\epsilon(t-z)\alpha_t(B)\,dt.
 \label{ha:gaussian-time}
\end{equation}
The second integral is entire between the relevant energy Sobolev
spaces, with norm loss at most $e^{(\operatorname{Im}z)^2/(2\epsilon^2)}$.
It agrees with complex conjugation on finite bands.  This supplies the
meaning of complex conjugation without applying $e^{sH}$ to an arbitrary
smooth-energy vector.  Moreover $B_\epsilon\to B$ strongly on
$\mathcal D_\infty$ in each fixed output energy norm.

Suppose a scalar spatial support and the support of $B$ have a strict
locality margin $d>0$, so $[a(g),\alpha_t(B)]=0$ for $|t|<d$ as forms.
Then only the Gaussian tail contributes, and for a sufficiently large
fixed $q'$,
\begin{equation}
 \|(1+H)^{-q'}[a(g),\alpha_{is}(B_\epsilon)](1+H)^{-q'}\|
 \leq Cp_s(g)\mathcal N(B)
          e^{s^2/(2\epsilon^2)}e^{-d^2/(4\epsilon^2)}.
 \label{ha:gaussian-tail}
\end{equation}
To obtain the time-zero locality premise, first smear in a small real-time
interval where spacelike separation holds; the all-weight bounds make
matrix elements smooth, and the temporal mollifier can be removed.
The form estimate \eqref{ha:scalar-form-bound}, rather than an ordinary
time-zero operator assertion, justifies both product orders.

Expand the commutator of \eqref{ha:stencil} with $B_\epsilon$ by the
product rule.  Each scalar factor is replaced by
$[a(g_a),\alpha_{is_a}(B_\epsilon)]$ and the original positive internal
gaps are unchanged.  This algebra is justified with spectral cutoffs
and a slightly larger form weight.  Indeed, if
$C=(1+H)^{-q}a(g)(1+H)^{-q}$ is bounded and $R=(1+H)^{-1}$, then
\begin{equation}
 \|R(C-P_ECP_E)R\|\leq\frac{2\|C\|}{1+E}.
 \label{ha:cutoff-form}
\end{equation}
Thus one extra weight on each side gives norm convergence, and fixed
positive gaps absorb all resulting finite polynomial weights.  There
is no compact-resolvent assumption.

Using \eqref{ha:gaussian-tail} at the replaced factor, summing the fixed
number of factors and their nuclear packet decomposition gives
\begin{equation}
 \|e^{-aH}[\Phi_G^{\rm av}(f),B_\epsilon]e^{-aH}\|
 \leq Ce^{CG}a^{-D_2}\mathcal N(f)\mathcal N(B)
       e^{\delta_G^2/(2\epsilon^2)}e^{-d^2/(C_2\epsilon^2)},
 \qquad a>2\rho_G.
 \label{ha:finite-commutator}
\end{equation}
Strict spacelike separation of compact center supports supplies a
uniform $d$ for all sufficiently small stencils.  Identity terms commute
exactly, and every lower-subset term obeys the same estimate.

Choose $\epsilon_G=a_G=c_0G^{-1/2}$ with $c_0$ small enough that the
Gaussian exponent in \eqref{ha:finite-commutator} dominates $CG$.
The finite-stencil commutator tends to zero exponentially.  To replace
it by the limit, use \eqref{ha:shrinking-heat-error} at $a_G/2$ and the
one-sided/all-weight bound for $B$.  The required crossing factor has
norm at most $Ca_G^{-p_B}e^{Ca_G^2/\epsilon_G^2}$, so the additional
loss is only a fixed power of $G$.  Increase the fixed native exponent
$K$ beyond those powers.  This gives
\begin{equation}
 \|e^{-a_GH}[\Phi(f),B_{\epsilon_G}]e^{-a_GH}\|\longrightarrow0.
 \label{ha:limiting-commutator}
\end{equation}
On $\mathcal D_\infty$, heat tends to the identity in every energy norm
and $B_{\epsilon_G}\to B$ in the required norms.  Therefore
$[\Phi(f),B]=0$ on this domain.

First take $B$ to be the original local scalar field and its adjoint.
This proves relative locality with the scalar sector.  Then take $B$
to be another new field or its companion.  The first application
supplies its locality with each scalar constituent, so the same
argument proves mutual locality of the new family.  It requires no
Haag duality or assertion that the scalar net already generates every
new operator.  The result is common-domain locality, conditional on
the stated scalar and copied inputs.

\paragraph{Provenance.}
\path{work/shrinking_scalar_stencil_locality_2026-10-18.txt}.

\section{What the conditional chain would establish, and what it does not}
\label{ha:identification}

If the source-attachment and simultaneous chronological hypotheses in
Section~\ref{ha:scope} are established for the actual scalar germ families,
and the scalar reconstruction inputs in Section~\ref{ha:scalar} are
retained on the same regulator subsequence, the argument produces a
finite family of closable, mutually local, tempered translation fields
on one scalar Hilbert space.  They have a common invariant
$\mathcal D_\infty$, polynomial energy bounds, actual adjoint companions,
centered vacuum expectations, and a positive-energy translation
representation.  Their one-insertion Euclidean continuation is the
constructed heat form.  This is a substantial conditional conclusion;
none of these functional-analytic deductions repairs its missing source
admission premise.

Several different remaining questions must not be collapsed into one.
First, Gaussian rank and the local inverse in
\eqref{ha:matrix-inverse} do not give a lower bound for the full-state
Gram matrix.  The complete mixed coefficient includes scalar output,
conditional covariance, and cross terms with historical contributions.
A positive universal vacuum coefficient or a positive conditional
Gaussian comparison can be canceled by the remaining mixed terms.
The two-colour attachment and the corrected power
\eqref{ha:packing-correct} are ingredients for controlling these terms,
not a completed relative lower bound.  Zero or dependent limiting
directions have not been excluded.

Second, shrinking translation periods prove continuous translations,
not rotations.  The matched frame is built from selected scalar stencils,
causal prefix subtractions, and terminal matrices.  Its limit has not
been identified with the standard
$\operatorname{Sym}^2(\Lambda^2\mathbb R^4)$ representation of
$O(4)$, nor with its Wick-rotated Lorentz tensor.  Adjoining actual
reflected companions proves closability without making that identification.
Finite-symmetry identities and a local Gaussian tensor label do not
replace continuous covariance of the full-state limit.

Third, scalar phase selection is not selection of a compatible phase
for an unconstructed full local field hierarchy.  Even when a family
of decomposable operator bounds passes almost everywhere to phases,
an averaged ultraviolet coefficient need not equal a fixed phase's
coefficient.  A phasewise matching theorem must retain all necessary
central tests, or a genuine decomposable operator-norm asymptotic.
A nonzero positive-energy derivative witness proves neither a uniform
mass gap nor uniqueness of the original mixed vacuum.

Finally, a full Yang--Mills identification requires a compatible joint
construction of gauge-invariant curvature differential polynomials,
their mixing and local identities, and the relevant nonabelian
short-distance laws on the same theory.  A scalar sector, twenty-one
Gaussian quadratic symbols, or one formal dimension-six cubic matching
coefficient is not that construction.  The fixed-$q$ refined parity
estimate recently obtained for cubic analysis concerns the physical
inner chart only.  It gives no all-age dimension-six transport theorem,
no marked full-history cubic source theorem, and no replacement for
the admission gap identified here.  Other gauge groups and a spectral
mass gap are additional scope requirements if that broader formulation
is intended.

The narrow outstanding source estimate is concrete: admit the
preparation-changed joined coarse slots, their retained endpoint frames,
and their complete averaging/current collars uniformly at a fixed
weak ceiling, and establish the improved first-read bound on its
simultaneous contours.  Repeating the all-age bookkeeping cannot prove
that estimate.  Once it is supplied, the conditional pressure, domain,
and locality arguments above can be invoked with their stated inputs;
the Gram and covariance questions still require their separate proofs.
